\chapter{Introduction}
\noindent\rule{\linewidth}{2pt}
This chapter provides a brief description of the Airline Reservation System, including its background, motivation, problem statement, and the overall organization of the project report.

\section{Background}

The airline industry plays a significant role in modern transportation by enabling efficient travel across cities, countries, and continents. With the increasing demand for air travel, airlines must manage large volumes of passenger data, flight schedules, and booking operations. Traditional manual reservation systems are often inefficient, time-consuming, and prone to human errors.

An Airline Reservation System is a software application designed to automate the process of flight booking, seat allocation, and passenger management. It allows users to search for flights, check seat availability, book tickets, and cancel reservations. The use of automated systems improves accuracy, reduces manual effort, and enhances customer satisfaction.

\section{Motivation for the present research work}

The motivation behind this project arises from the limitations of manual and semi-automated reservation systems. These systems often lead to booking conflicts, incorrect passenger records, and inefficient seat allocation. As the number of air travelers continues to increase, there is a growing need for reliable and efficient reservation systems.

This project aims to develop an Airline Reservation System using Object-Oriented Programming (OOP) in Java. The goal is to design a system that is modular, scalable, and easy to maintain, while providing accurate and fast booking services.

\section{Problem statement}

In the current digital era, managing airline reservations efficiently is a significant challenge for many organizations. Existing systems may lack proper structure, scalability, and data consistency, leading to issues such as double booking, incorrect data handling, and delays in processing.

The problem addressed in this project is to design and develop an automated Airline Reservation System that ensures accurate booking management, efficient seat allocation, and secure handling of passenger data.

\section{Organization of the project report}

The research work presented in this report is organized into several chapters:

\begin{itemize}
\item \textbf{Chapter 1} describes the introduction, including background, motivation, and problem statement.
\item \textbf{Chapter 2} presents the literature review of existing systems.
\item \textbf{Chapter 3} explains the methodology used in the development of the system.
\item \textbf{Chapter 4} describes the system architecture.
\item \textbf{Chapter 5} discusses the implementation details.
\item \textbf{Chapter 6} presents results and analysis.
\item \textbf{Chapter 7} concludes the project and suggests future enhancements.
\end{itemize}